\section{Cuencas muestreadas y transporte de las condiciones iniciales}
\label{bre:bas:section}

Los retratos de fase describen la trayectoria alcanzada desde una semilla;
los sondeos y cortes siguientes muestran cómo cambia el destino al cambiar
el dato inicial. En Caputo, cada punto del muestreo define un PVI nuevo con
terminal inferior cero y la misma convención de reinicio. Los colores
clasifican trayectorias en una ventana finita. La continuación, en cambio,
transporta una condición inicial entre sistemas: su gráfica debe leerse
junto con el parámetro que cambia y el tratamiento de la memoria.

\subsection{Chua no suave: sondeos dentro de bolas}
\label{bre:bas:pwl}

Para el caso de la Sec.~\ref{bre:res:chua-pwl}, con $q=0.9998$, se
muestrearon los tres equilibrios mediante
$X_0=e+rU^{1/3}d$, con $U$ uniforme en $[0,1]$ y $d$ uniforme en la
esfera unidad, independientes. Por tanto, $r$ es el radio de la bola y
la distancia de una sonda concreta puede ser menor que $r$. Se usaron
ABM--PECE con memoria completa, $h=0.01$, horizonte máximo $T=200$ y
descarte de 50. El criterio de contacto fue el percentil 90 de distancias
euclidianas a la nube objetivo, con umbral $0.5$ y hasta 2000 puntos de
cada cola. Los criterios de equilibrio y salida pudieron detener antes
una integración; una salida numérica no demuestra explosión matemática.

El protocolo local reunió 7200 trayectorias en los siete radios desde
$10^{-5}$ hasta $10^{-2}$ y no registró contactos. La extensión a
$r=0.03,0.1,0.3$ añadió 10\,200 trayectorias. Hasta $r=0.1$ el
acumulado fue $0/13\,800$; en las 3600 sondas de $r=0.3$ aparecieron
37 contactos, 19 desde la bola de $E_+$ y 18 desde la de $E_-$.
Sus distancias iniciales efectivas estuvieron entre $0.1941805298$ y
$0.2997254735$. El primer radio ensayado con acceso delimita una escala
de búsqueda; no estima la distancia mínima exacta a la cuenca.

\begin{figure*}[t]
 \centering
 \includegraphics[width=0.74\textwidth]{pwl_sondeos_volumetricos.pdf}
 \caption{Chua no suave de Caputo: nueve historias representativas de un
 banco de 17\,400 sondas volumétricas. Se distinguen el objetivo, los
 equilibrios y las trayectorias con destinos diferentes. Los triángulos
 señalan el primer cruce del encuadre por las dos historias de salida
 representadas; no son estados finales. La figura no dibuja la frontera
 de una cuenca ni representa todas las sondas.}
 \label{bre:bas:fig-pwl-volume}
\end{figure*}

La Fig.~\ref{bre:bas:fig-pwl-volume} permite comparar transitorios que
parten de vecindades semejantes y terminan en clases distintas. Sus ejes
conservan las coordenadas físicas, con límites ajustados por componente;
la caja tridimensional no impone una misma escala numérica a los tres
ejes. Los contactos a radio macroscópico se interpretan separadamente de
la ausencia observada dentro del protocolo local.

\begin{figure*}[t]
 \centering
 \includegraphics[width=0.98\textwidth]{pwl_nube_volumetrica_conteos.pdf}
 \caption{Condiciones iniciales y destinos del sondeo volumétrico PWL.
 Arriba: las 1200 condiciones efectivamente integradas en la bola de
 radio $0.3$ de cada equilibrio, trasladadas a coordenadas locales
 $X_0-e$, con la misma escala física en los tres ejes. El signo $+$
 marca el centro; no se normalizan ni se superponen bolas de radios
 diferentes. Abajo: conteos apilados de los diez radios, 5800
 trayectorias por equilibrio. Los colores indican la clasificación de
 horizonte finito obtenida mediante ABM--PECE con $h=0.01$, memoria
 completa, $T\leq200$ y descarte de 50. El amarillo muestra los
 37 contactos, todos en las bolas de radio $0.3$ de $E_\pm$.
 La densidad visual del panel superior no representa una densidad
 invariante del flujo ni una probabilidad certificada de atracción.}
 \label{bre:bas:fig-pwl-clouds}
\end{figure*}

La Fig.~\ref{bre:bas:fig-pwl-clouds} representa las condiciones iniciales
para el radio mayor ensayado. Los equilibrios de
\eqref{bre:res:chua-equilibrium-reduction} usados en la traslación son
\[
 E_0=0,\qquad E_\pm\simeq\pm\begin{pmatrix}
 6.5883078865\\0.0028364023\\-6.5854714843
 \end{pmatrix}.
\]
Todas las condiciones iniciales satisfacen $\norm{X_0-e}/r\leq1$.
Los conteos suman 5827 destinos de equilibrio, 5826 de otra clase,
5710 salidas y 37 contactos.

\subsection{Chua arctangente: superficies esféricas y cortes planos}
\label{bre:bas:arctan}

En c590 se mantienen los parámetros de
\eqref{bre:res:c590-parameters} y $q=0.9999$. Las 13\,500 sondas
pertenecen a superficies $X_0=e+rd$, mediante direcciones axiales y una
distribución de Fibonacci sobre la esfera. Se usaron nueve radios,
$10^{-5},10^{-4},10^{-3},10^{-2},0.03,0.1,0.3,1,2$, con
100, 200, $\ldots$, 900 direcciones por equilibrio, respectivamente.
No se trata de un muestreo uniforme del volumen de las bolas.
Las integraciones utilizaron ABM con memoria completa, $h=0.005$,
$T=180$ y descarte de 90, conservando las reglas de parada del protocolo.

Hasta $r=0.3$ hubo cero contactos entre 8400 trayectorias. En $r=1$
hubo 22 de 2400, todos desde $E_0$; en $r=2$, 588 de 2700, repartidos
en 289, 150 y 149 desde $E_0,E_+,E_-$. El total fue 610 contactos.
Los otros destinos fueron 4 equilibrios, 131 salidas y 12\,755
trayectorias de otra clase. Esta última categoría no determina por sí
sola un atractor diferente: puede contener transitorios todavía no
resueltos por la ventana y el clasificador.

\begin{figure*}[t]
 \centering
 \includegraphics[width=0.65\textwidth]{arctan_sondeos_superficiales.pdf}
 \caption{c590: nueve historias representativas seleccionadas de los
 sondeos sobre superficies esféricas hasta $r=1$. El conteo completo de
 13\,500 pruebas incorpora además el bloque $r=2$, que no se representa
 mediante historias adicionales en esta figura. Las circunferencias
 aparentes son proyecciones de trayectorias; no delimitan bolas de muestreo
 ni prueban periodicidad exacta del PVI de Caputo.}
 \label{bre:bas:fig-arctan-sphere}
\end{figure*}

Un análisis complementario fijó $z(0)$ y clasificó datos
$(x(0),y(0),z(0))$ en cinco planos. Los cortes se eligieron a partir de
las coordenadas de los equilibrios y los percentiles 10, 50 y 90 de la
nube postransitoria simétrica $A_+\cup A_-$:
\[
 z(0)\in\{-12.56295,-4.64761,0,4.64761,12.56295\}.
\]
El rectángulo común fue
$x(0)\in[-11.238014,11.238014]$,
$y(0)\in[-4.311997,4.311997]$. Se mantuvieron $h=0.005$, $T=180$
y descarte de 90. El contacto con una nube de referencia se evaluó
mediante el percentil 90 y umbral $0.5901016489$; la tolerancia de
clasificación a equilibrio fue $0.5$. Estas tolerancias describen
vecindad numérica en coordenadas físicas, no igualdad de conjuntos límite.

\begin{figure*}[t]
 \centering
 \includegraphics[width=\textwidth]{c590_cuencas_multicorte.pdf}
 \caption{Clasificación por destinos de c590 sobre cinco cortes $xy$.
 Las estrellas identifican equilibrios contenidos en el plano; los
 círculos sólo proyectan equilibrios de otros planos, con su separación
 $\Delta z$. La retícula visual es $241\times241$ por corte y contiene
 relleno categórico: los 290\,405 nodos mostrados corresponden a
 62\,416 trayectorias efectivamente integradas, no a una integración
 independiente por píxel. El gris incluye relleno y omisiones heredadas;
 la etiqueta gráfica de divergencia sólo se interpreta como resultado
 dinámico en los nodos cuya integración fue realizada.}
 \label{bre:bas:fig-c590-slices}
\end{figure*}

La retícula se refinó de forma adaptativa: se conservaron 35\,941
integraciones de la malla precedente y se integraron 26\,475 centros
de celdas cuyos cuatro vértices tenían clases distintas. Los nodos
restantes recibieron relleno categórico sin suavizado. Una máscara
separa los nodos calculados de ese relleno y de las omisiones grises
heredadas; por ello las áreas coloreadas no proporcionan directamente
fracciones de volumen de las cuencas. Entre los nodos realmente
integrados hubo 17\,369 contactos con $A_+$, 17\,514 con $A_-$,
843 destinos de equilibrio, 10\,331 de otra clase acotada y 16\,359
salidas, sin fallos numéricos registrados. Los cortes muestran regiones
de acceso y zonas donde conviene refinar; no establecen conectividad
global, dimensión de frontera ni exclusión de una vecindad tridimensional.

\subsection{Chua entero: destinos cerca de los equilibrios}
\label{bre:bas:integer}

Para el Chua entero de la Sec.~\ref{bre:res:chua-integer}, el control
comparado tomó las mismas 504 condiciones en las dos implementaciones:
tres equilibrios, siete radios entre $10^{-5}$ y $10^{-2}$ y 24 puntos
uniformes dentro de cada bola. Con $T=500$, descarte de 120,
tolerancia de equilibrio $10^{-3}$ y tolerancia normalizada a la nube
$0.08$, coincidieron las 504 clases: 152 trayectorias hacia $E_0$,
198 acotadas no triviales y 154 salidas. Ninguna hizo contacto con el
candidato. La concordancia reproduce el resultado del muestreo común;
no amplía sus radios ni convierte sus etiquetas finitas en una prueba
asintótica.

\begin{figure*}[t]
 \centering
 \includegraphics[width=0.72\textwidth]{chua_entero_sondeos.png}
 \caption{Trayectorias del control de vecindades del Chua entero, con la
 nube objetivo en verde. Los ejes $x0,x1,x2$ designan las tres componentes
 del estado $(x,y,z)$, no tiempos iniciales. La escala amplia conserva
 las trayectorias que abandonan la región del candidato; el retrato
 detallado del objetivo aparece en la Fig.~\ref{bre:res:fig-chua-integer}.}
 \label{bre:bas:fig-integer-probes}
\end{figure*}

\subsection{MAVPD: selección por continuación y sondeo del candidato}
\label{bre:bas:mavpd}

La continuación descrita en la Sec.~\ref{bre:res:mavpd} transportó la
rama de menor frecuencia desde $(\xi,\gamma)=(3.1,0.1)$ hasta
$\xi=2.85$, y después modificó $\gamma$. La
Fig.~\ref{bre:bas:fig-mavpd} compara el exponente máximo finito con
$-\max\operatorname{Re}\lambda(DF(E_\pm))$, positivo cuando los
equilibrios laterales son linealmente estables. El exponente de una
trayectoria y el margen local de un equilibrio miden objetos distintos.
Los siete desplazamientos respecto de
$\gamma_H=0.14380379839949112$ no definen una curva de bifurcación
continua certificada: los segmentos sólo enlazan valores ensayados.

En los desplazamientos $0.002,0.003,0.005$, el sondeo preliminar de
$E_0$ dio 6 contactos de 12. En $0.008$ y $0.010$ hubo cero contactos
y exponentes máximos mayores que $0.5$. La regla prefijada seleccionó
el mayor de estos desplazamientos, $\gamma=0.15380379839949113$;
un exponente positivo aislado no fue el único criterio.
En cada nivel del transporte se descartaron 60 unidades y se retuvieron
10, con DOP853, tolerancias $2\times10^{-10}$ y $2\times10^{-12}$,
y paso máximo $0.02$.

\begin{figure*}[t]
 \centering
 \includegraphics[width=0.49\textwidth]{mavpd_continuacion_cribado.pdf}\hfill
 \includegraphics[width=0.49\textwidth]{mavpd_sondeos_destinos.pdf}
 \caption{MAVPD. Izquierda: cribado de los niveles de continuación en
 $\gamma$; la tasa azul es un exponente de tiempo finito y la naranja
 es el margen lineal de $E_\pm$. Derecha: destinos de las 112 sondas
 del candidato seleccionado. Las barras 40, 36 y 36 cuentan condiciones
 lanzadas alrededor de $E_0,E_+,E_-$ que terminaron en algún equilibrio;
 no indican que todas hayan vuelto al equilibrio del que partieron.
 La barra de contactos con el objetivo es nula en los tres grupos.}
 \label{bre:bas:fig-mavpd}
\end{figure*}

El sondeo final utilizó 108 direcciones esféricas deterministas: tres
equilibrios, radios $10^{-5},10^{-3},10^{-2}$ y doce direcciones por
radio. Se añadieron cuatro condiciones en ambos sentidos del autovector
inestable de $E_0$, a distancias $10^{-7}$ y $10^{-4}$. Se descartaron
400 unidades y se conservaron 100; DOP853 usó tolerancias
$2\times10^{-10}$ y $2\times10^{-12}$ y paso máximo $0.03$.
Todas las sondas terminaron en un equilibrio, sin contactos con el
candidato, ambigüedades ni fallos. Es una exploración localizada de
destinos cerca de los equilibrios, no un mapa volumétrico completo.

\subsection{PLL: continuación de la ganancia y vecindades cilíndricas}
\label{bre:bas:pll}

Con $(\tau_1,\tau_2,\omega_\Delta)=(0.0448,0.0185,178.9)$, la
continuación aumenta la ganancia $L$ desde cero hasta 500 en pasos de 5.
La Fig.~\ref{bre:bas:fig-pll} representa la velocidad angular y la
coordenada del filtro en la sección $\theta=0$. Cada punto representa
un ciclo obtenido del problema de retorno del nivel correspondiente;
no es una muestra consecutiva de una trayectoria a parámetros fijos.
En $L=0$, la coordenada $x$ procede de la solución explícita de la
Sec.~\ref{bre:res:pll}, pues la inversión basada en dividir por $L$
sólo es válida cuando $L>0$.

\begin{figure*}[t]
 \centering
 \includegraphics[width=0.94\textwidth]{pll_continuacion_ganancia.pdf}\\[4pt]
 \includegraphics[width=0.66\textwidth]{pll_sondeos_cilindricos.pdf}
 \caption{PLL. Arriba: continuación de las coordenadas de retorno hasta
 $L=500$. Abajo: las 96 condiciones iniciales de los sondeos, expresadas
 mediante $u=x-x_e$ y la fase desplazada y envuelta $v$. Los dos grupos
 corresponden al foco y a la silla. La figura inferior representa datos
 iniciales, no una frontera de cuenca; sus 96 trayectorias terminaron en
 el foco bajo el protocolo indicado.}
 \label{bre:bas:fig-pll}
\end{figure*}

En el nivel final se usó la distancia cilíndrica escalada con escalas
$s_u=\tau_1/2=0.0224$ y $s_v=\pi$. Alrededor del foco y de la silla
se tomaron doce puntos uniformes dentro de cada bola escalada, para
$r=10^{-5},10^{-4},10^{-3},10^{-2}$: 96 condiciones en total.
El horizonte fue 5 y la cola evaluada tuvo longitud 1, con paso máximo
$0.002$, tolerancia al foco $10^{-5}$ y umbral de contacto con la nube
$0.01$. Las 48 sondas de cada equilibrio terminaron en el foco.
La coexistencia del ciclo corredor con el foco y la separatriz calculada
explica la utilidad de continuar un ciclo lejano; la exclusión matemática
de vecindades completas exigiría validar la separación de cuencas.

\subsection{Kalman--Fitts: del retorno por conmutación al ciclo suave}
\label{bre:bas:kalman}

La Fig.~\ref{bre:bas:fig-kalman} completa la ruta de la
Sec.~\ref{bre:res:kalman}: a partir del retorno del sistema con signo se
transporta el estado por los 51 niveles de
$(1-\lambda)\operatorname{sign}(c^{\T}X)+
\lambda\tanh(c^{\T}X/0.01)$. Los puntos superiores son estados al
finalizar cada etapa, de duración 120 y paso $0.01$; su variación refleja
tanto el cambio de sistema como la fase alcanzada. No son amplitudes
extremas, puntos fijos ni un diagrama de bifurcación.

\begin{figure*}[t]
 \centering
 \includegraphics[width=0.78\textwidth]{kalman_continuacion_signo_tanh.pdf}\\[4pt]
 \includegraphics[width=0.96\textwidth]{kalman_ciclo_proyecciones.pdf}
 \caption{Kalman--Fitts. Arriba: los cuatro estados terminales de la
 continuación desde signo hasta tangente hiperbólica. Abajo: proyecciones
 del régimen final del sistema suave, con $(x,y,z)=(x_1,x_2,x_3)$;
 la cuarta coordenada permanece en la integración. La geometría se
 contrasta con el periodo numérico $11.9169960234$ y los controles
 de retorno e integración de la Sec.~\ref{bre:res:kalman}.}
 \label{bre:bas:fig-kalman}
\end{figure*}

El control local alrededor del único equilibrio utilizó 48 datos
uniformes dentro de bolas de cuatro dimensiones: doce para cada radio
$10^{-5},10^{-4},10^{-3},10^{-2}$. Con horizonte 600, descarte 150 y
paso $0.02$, ninguno contactó con el ciclo. Cuarenta recibieron etiqueta
de equilibrio y ocho quedaron acotados sin contacto bajo la tolerancia
utilizada. Estos ocho resultados no se reinterpretan como convergencias
asintóticas ya resueltas. La trayectoria final, los retornos y el sondeo
local aportan controles complementarios a la continuación.
